\section{Keywords in Patterns: \texttt{my-if2}}
\label{sec:literals}

Every \code{syntax-case} so far had an empty list \code{()} after \code{stx}. That list names \emph{literals}: identifiers that a pattern matches exactly, instead of binding them as pattern variables. They let a macro have keywords of its own:
\begin{lstlisting}
(define-syntax (my-if2 stx)
  (syntax-case stx (then else)
    [(my-if2 e1 then e2 else e3)
     #'(if e1 e2 e3)]))
\end{lstlisting}
In the pattern, \code{e1}, \code{e2}, and \code{e3} are pattern variables, but \code{then} and \code{else} are listed as literals, so the use must contain those words in those positions:
\begin{lstlisting}
(my-if2 (> 2 1) then "yes" else "no")    ; expands to (if (> 2 1) "yes" "no")
\end{lstlisting}
A use without the keywords, such as \code{(my-if2 (> 2 1) "yes" "no")}, matches no clause, and the expander reports a syntax error for \code{my-if2} before the program runs.

\begin{pitfall}[Forgetting the literal list]
Without \code{(then else)}, the pattern \code{(my-if2 e1 then e2 else e3)} would treat \code{then} and \code{else} as two more pattern variables. Every use with five subforms would match, whatever words stood in those positions, and \code{(my-if2 c 1 2 3 4)} would expand to \code{(if c 2 4)} without complaint.
\end{pitfall}

\begin{remark}
A literal is matched as an identifier, by binding, not only by spelling: the \code{else} in the use must refer to the same \code{else} as the one in the macro's literal list. This is hygiene again (\Cref{sec:hygiene}), applied to keywords. \code{cond}'s \code{else} works the same way.
\end{remark}
